\chapter{基于时空方差导向（SVGF）的降噪算法实现}

\section{极低采样率（1 SPP）蒙特卡洛积分的噪声分析}

渲染方程本质上是对半球空间连续函数的无穷积分。在物理真实的路径追踪算法中，通常使用蒙特卡洛方法进行离散采样近似：

\begin{equation}
L_o \approx \frac{1}{N} \sum_{i=1}^{N} \frac{f_r(\omega_i, \omega_o) L_i(\omega_i) (n \cdot \omega_i)}{p(\omega_i)}
\end{equation}

其中 $N$ 为采样射线数量，$p(\omega_i)$ 为采样的概率密度函数（PDF）。根据大数定律，当 $N \to \infty$ 时，近似解收敛于真实值。然而，在实时渲染中 $N=1$。这种极度欠采样的状态导致估计值偏离期望，产生高斯分布特性的方差（即噪点）。

传统的双边滤波（Bilateral Filter）等纯空域降噪算法在处理此类高方差信号时，极易导致几何边缘模糊或纹理细节丢失（即“过度涂抹”现象）。因此，本系统引入了时空联合降噪策略，以期在保留高频细节的同时彻底抹平低频噪声。

\section{几何运动矢量（Motion Vector）与时域重投影}

时域滤波（Temporal Filtering）的核心思想是“向过去借用样本”。既然当前帧只有 1 SPP，如果能将过去数十帧的样本进行累积，即可在不增加单帧算力的前提下，等效获得几十 SPP 的高质量画面。其前提是必须精准找到当前帧像素在上一帧中的物理对应位置。

\subsection{运动矢量的精确计算}

在 G-Buffer 生成阶段，本系统为场景中的每个顶点不仅传入了当前帧的模型视图投影矩阵（MVP），还传入了上一帧的矩阵（Previous MVP）。在顶点着色器中，分别计算当前帧的裁剪空间坐标 $P_{clip}$ 与上一帧的裁剪空间坐标 $P_{prev\_clip}$。

经过透视除法转换至 NDC 空间后，运动矢量 $V_{motion}$ 可定义为：

\begin{equation}
V_{motion} = P_{ndc}^{current} - P_{ndc}^{prev}
\end{equation}

该二维矢量被写入格式为 \texttt{VK\_FORMAT\_R16G16\_SFLOAT} 的 G-Buffer 纹理中，作为时域重投影的寻址偏移量。

\subsection{时域历史累积公式}

在 SVGF 的时域处理阶段（Temporal Reprojection Pass），系统通过当前像素的 UV 坐标减去 $V_{motion}$，采样历史帧的光照缓冲与方差缓冲。对于光照颜色 $C$，其指数滑动平均（EMA）累积公式为：

\begin{equation}
C_{current} = \alpha C_{raw} + (1 - \alpha) C_{history}
\end{equation}

其中 $C_{raw}$ 为当前帧光追输出的带噪点颜色，$C_{history}$ 为重投影获取的历史颜色。$\alpha$ 为混合权重（通常取值在 0.1 到 0.2 之间），决定了历史数据的留存比例。

\section{局部方差实时估计与历史有效性验证}

时域累积存在一个致命缺陷：当场景发生剧烈光照变化（如光源突然移动）或物体相互遮挡（Occlusion）时，由于历史像素的物理状态已失效，强行累积会导致严重的视觉拖影（Ghosting）伪影。

\subsection{历史有效性剔除（History Rejection）}

为了消除拖影，系统在重投影后，必须对获取的历史像素进行有效性验证。本系统采用深度与法线联合一致性检验：

\begin{enumerate}
    \item \textbf{深度测试}：比较当前像素的线性深度 $z_{current}$ 与历史像素的线性深度 $z_{history}$，如果相对误差超过阈值，则判定发生了遮挡断层。
    \item \textbf{法线测试}：计算当前法线 $N_{current}$ 与历史法线 $N_{history}$ 的点积，若点积小于设定的角度余弦阈值（如 0.9），则判定几何表面发生突变。
\end{enumerate}

当任一验证失败时，系统将动态增大混合权重 $\alpha$（甚至强制令 $\alpha = 1.0$），即彻底抛弃历史数据，完全依赖当前帧的新信号。

\subsection{空间方差动态估计}

在时域断层处（如物体移动边缘），由于历史数据被丢弃，噪点会瞬间爆发。此时系统必须动态计算局部空间方差（Spatial Variance），以引导后续的空域滤波强度。系统在一个 $7 \times 7$ 的局部邻域窗口内，通过计算一阶矩（亮度均值 $\mu$）与二阶矩（亮度平方均值 $\mu^2$）来实时估计方差 $\sigma^2$：

\begin{equation}
\sigma^2 = \max(0, \mu^2 - (\mu)^2)
\end{equation}

该方差图与光照图同步输出，作为下一阶段空域滤波的关键输入。

\section{基于 A-Trous 小波变换的边缘保留空域滤波}

为了进一步处理时域累积无法消除的残余噪声（尤其是在历史数据被丢弃的区域），本系统在管线末端引入了多级 A-Trous（带孔）小波滤波算法。

\subsection{A-Trous 扩张卷积架构}

传统的高斯模糊若要获得较大的感受野，需要极大的卷积核，导致 GPU 性能急剧下降。A-Trous 算法通过在每一轮迭代（Iteration）中对固定大小的 $3 \times 3$ 卷积核插入“孔洞”（即将采样步长 Step 翻倍），使得感受野以 $2^i$ 的指数级急剧扩大。仅需 4 至 5 次迭代，即可覆盖极大的屏幕像素范围，且每次迭代的计算时间保持常数级。

\subsection{联合边缘停止函数（Edge-Stopping Functions）}

为了防止跨越几何边界导致画面糊化，A-Trous 滤波器在计算邻域像素 $q$ 对中心像素 $p$ 的权重贡献 $W_{pq}$ 时，联合了 G-Buffer 提供的多维几何信息，构建了严苛的边缘停止函数：

\begin{equation}
W_{pq} = w_z \cdot w_n \cdot w_l
\end{equation}

\begin{itemize}
    \item \textbf{深度权重 $w_z$}：防止跨越前景与背景的深度断层。考虑到屏幕空间的深度梯度，引入偏导数项 $\nabla z$ 进行自适应放宽。
    \item \textbf{法线权重 $w_n$}：防止跨越物体折角与边缘：
    \begin{equation}
    w_n = \max(0, n_p \cdot n_q)^{\sigma_n}
    \end{equation}
    \item \textbf{亮度方差权重 $w_l$}：由 5.3 节计算的局部方差 $\sigma^2$ 直接引导。方差越大的高噪点区域，滤波权重越大（允许更强烈的模糊）；方差越小的平滑区域，滤波权重越小（保护高频细节）。
\end{itemize}

通过这三项权重的严格限制，系统成功实现了“在平滑区域强力抹平噪点，在几何边缘与纹理细节处精准保留”的智能滤波效果。

\section{本章小结}

本章详细探讨了混合渲染管线中不可或缺的 SVGF 降噪子系统的架构与实现。针对 1 SPP 蒙特卡洛积分产生的剧烈方差，系统设计了时空解耦的联合降噪策略。首先，利用 G-Buffer 提取的精确运动矢量，实现了基于指数滑动平均的时域历史累积，并引入深度与法线一致性检验，有效解决了动态场景下的拖影伪影问题。其次，针对历史数据失效区域，系统实时估算局部空间方差，并以此引导多轮 A-Trous 扩张空域滤波，结合法线、深度与亮度的多维边缘停止函数，实现了高保真度的抗噪平滑。SVGF 算法的成功集成，使得本渲染引擎能够在主流消费级显卡上，以实时的性能开销提供接近离线光线追踪的物理级视觉画质。
